% Causal Knowledge Base Editor poster — editorial layout (A0 portrait).
% Reuses the AgonProject poster style: the style layer lives in
% agoneditorial.sty (colours, title/background styles, and the helper macros
% \eyebrow, \screenshot, \featureitem, \shotpanel, \semrow, \callout,
% \sectionrule, \formalism). This file is a scaffold — content is still WIP.

\documentclass[25pt,a0paper,portrait,margin=0mm,innermargin=45mm,blockverticalspace=10mm,colspace=18mm,subcolspace=10mm]{tikzposter}

\usepackage[T1]{fontenc}
\usepackage[utf8]{inputenc}
\usepackage{sourceserifpro}
\usepackage{sourcesanspro}
\usepackage{sourcecodepro}
\usepackage{booktabs}
\usepackage{graphicx}
\usepackage{qrcode}
\usepackage{fontawesome5}
\usepackage{xspace}

\usepackage{agoneditorial}

\newcommand{\ckbe}{Causal Knowledge Base Editor\xspace}
\newcommand{\tweety}{TweetyProject\xspace}

\usecolorstyle{editorialcolor}
\usebackgroundstyle{editorialback}
\usetitlestyle{editorialtitle}
\useblockstyle{flat}

\title{Argumentation-based Causal Reasoning}
\subtitle{Lars Bengel \textbullet\ Oleksandr Dzhychko \textbullet\ Myriam Kasten \textbullet\ Kai Sauerwald \textbullet\ Matthias Thimm}

\begin{document}
\maketitle

%======================================================================
% Lede
%======================================================================
\begin{columns}
\column{0.56}
\block{}{
    \begin{minipage}[c][0.6\linewidth][c]{\linewidth}
    \eyebrow{Web tool \textbullet\ open source}
    \vspace{0.5em}

    {\sffamily\bfseries\color{ink}\fontsize{66}{70}\selectfont
        Build causal models and reason about \textcolor{aigblue}{why}.\par}
    \vspace{0.7em}

    {\color{inksoft}\fontsize{27}{35}\selectfont
        The {\sc\ckbe} is a visual interface to construct causal knowledge bases,
        \textbf{reason} over them through their induced argumentation framework,
        and get \textbf{human-readable explanations}.\par}
    \vspace{1.3em}

    {\centering\begin{tcolorbox}[hbox, colback=aigyellow, colframe=aigyellow, boxrule=0pt, arc=14pt,
        left=18pt, right=22pt, top=16pt, bottom=16pt, boxsep=0pt]
        \begin{minipage}[c]{3.9cm}
            \centering
            \setlength{\fboxsep}{8pt}\colorbox{white}{\qrcode[height=3.4cm]{https://causal-knowledge-base-editor.aig.fernuni-hagen.de}}
        \end{minipage}\hspace{22pt}%
        \begin{tabular}[c]{@{}l@{}}
            {\sffamily\bfseries\color{aigbluedeep}\fontsize{24}{28}\selectfont\textls[150]{SCAN TO TRY IT LIVE}}\\[0.35em]
            {\sffamily\bfseries\color{aigbluedeep}\fontsize{30}{36}\selectfont causal-knowledge-base-editor.aig.fernuni-hagen.de}
        \end{tabular}%
    \end{tcolorbox}\par}
    \end{minipage}
}
\column{0.44}
\block{}{
    \vspace{1.2em}
    \centering\screenshot[0.96\linewidth]{figures/model.pdf}\par
    \vspace{0.6em}
    {\color{muted}\itshape\fontsize{20}{25}\selectfont Graphical editor --- a causal model
     for diagnosing a viral disease, with structural equations and assumptions.\par}
}
\end{columns}

%======================================================================
% How it works — the argumentation-based causal reasoning pipeline
%======================================================================
\begin{columns}\column{1}\block{}{\sectionrule\vspace{0.6em}\eyebrow{How it works}}\end{columns}
\begin{columns}
\column{0.31}\block{}{\featureitem{1 \enspace Causal model}{\faSitemap}{%
    Boolean-valued variables tied together by \emph{structural equations}
    $v \leftrightarrow \phi$ --- one causal mechanism per explainable atom --- plus
    assumptions on the uncontrollable background atoms.}}
\column{0.31}\block{}{\featureitem{2 \enspace Induced AF}{\faProjectDiagram}{%
    Arguments are minimal sets of assumptions that consistently entail a conclusion;
    an argument \emph{undercuts} another when it negates one of its premises,
    yielding an argumentation framework.}}
\column{0.31}\block{}{\featureitem{3 \enspace Reason \& explain}{\faBalanceScale}{%
    Given an observation, a conclusion holds exactly when it is in \emph{every stable
    extension} of the induced AF --- with dependency and sequence-based explanations.}}
\end{columns}

%======================================================================
% Features
%======================================================================
\begin{columns}\column{1}\block{}{\sectionrule\vspace{0.6em}\eyebrow{Features}}\end{columns}
\begin{columns}
\column{0.2}\block{}{\featureitem{Model}{\faPenNib}{Construct causal models graphically --- atoms, structural equations, automatic layout.}}
\column{0.2}\block{}{\featureitem{Assume \& observe}{\faSlidersH}{Fix assumptions on background atoms and choose the observations to reason from.}}
\column{0.2}\block{}{\featureitem{Reason}{\faBalanceScale}{Valid conclusions computed from the stable extensions of the induced AF.}}
\column{0.2}\block{}{\featureitem{Explain}{\faStream}{Dependency explanations and sequence-based explanations of conclusion arguments.}}
\column{0.2}\block{}{\featureitem{Switch views}{\faRetweet}{Flip freely between the causal model and the induced AF behind it.}}
\end{columns}

%======================================================================
% In action
%======================================================================
\begin{columns}\column{1}\block{}{\sectionrule\vspace{0.6em}\eyebrow{In action}}\end{columns}
\begin{columns}
\column{0.5}\block{}{\shotpanel[\linewidth]{Causal model}{figures/model.pdf}{A causal model for diagnosing a viral disease --- explainable atoms, their structural equations and the background assumptions.}}
\column{0.5}\block{}{\shotpanel[\linewidth]{Induced AF \& evaluation}{figures/af.pdf}{The argumentation framework induced from the model, with the valid conclusions read off its stable extensions.}}
\end{columns}
\begin{columns}
\column{0.5}\block{}{\shotplaceholder{Assumptions \& observations}{The input panel where the user fixes background assumptions and observations before evaluating. \emph{Screenshot to provide.}}}
\column{0.5}\block{}{\shotplaceholder{Sequence-based explanation}{The step-by-step explanation of a conclusion argument. \emph{Screenshot to provide.}}}
\end{columns}

%======================================================================
% Outlook
%======================================================================
\begin{columns}
\column{0.58}
\block{}{
    \sectionrule\vspace{0.6em}
    \eyebrow{Outlook}
    \vspace{0.7em}
    {\color{inksoft}\fontsize{25}{33}\selectfont
    \begin{itemize}\setlength{\itemsep}{0.4em}
        \item richer presentation and explanation of results within the AF
        \item counterfactual reasoning via the twin-model method
    \end{itemize}}
}
\column{0.42}
\block{}{
    \sectionrule\vspace{0.6em}
    \eyebrow{Under the hood}
    \vspace{0.7em}
    {\color{inksoft}\fontsize{23}{30}\selectfont
     A JavaScript front end built on the open-source \texttt{aig-graph-component};
     causal reasoning and AF semantics by \emph{\tweety}.\par}
    \vspace{0.6em}
    \callout{Open source --- graphical modelling made to be explored, extended and reused.}
}
\end{columns}

%======================================================================
% Footer band — full-bleed, pinned to the bottom edge of the page
%======================================================================
\footerband{%
    {\ttfamily\bfseries\color{white}\fontsize{26}{30}\selectfont
     causal-knowledge-base-editor.aig.fernuni-hagen.de
     \quad\textbullet\quad \faGithub\ github.com/aig-hagen/aig-causal-knowledge-base-editor\par}
    \vspace{0.4em}
    {\ttfamily\color{bluetint}\fontsize{18}{23}\selectfont Lars Bengel \textbullet\ Oleksandr Dzhychko
     \textbullet\ Myriam Kasten \textbullet\ Kai Sauerwald \textbullet\ Matthias Thimm\par}%
}

\end{document}
